\section{Topologie metrischer Räume}
\label{sec:2.2}

\begin{definition}{Metrik, metrischer Raum}
    Sei \(X\) eine Menge. Eine Abbildung \(d : X \times X \to \mathbb{R}\) heißt Metrik, wenn für alle \(x,y,z \in X\) gilt:
    \[
        \text{(M1)}\; d(x,y)=0 \Rightarrow x=y,\qquad \text{(M2)}\; d(x,y)=d(y,x),\qquad \text{(M3)}\; d(x,z) \le d(x,y) + d(y,z).
    \]
    Ein metrischer Raum ist ein Paar \((X,d)\), bestehend aus einer Menge \(X\) und einer Metrik \(d\).
\end{definition}

\begin{definition}{Norm, normierter Raum}
    Sei \(V\) ein Vektorraum über \(\mathbb{R}\) oder \(\mathbb{C}\). Eine Abbildung \(\|\cdot\| : V \to \mathbb{R}\) heißt Norm, wenn für alle \(x,y \in V\) und \(\alpha\) gilt:
    \[
        \text{(N1)}\; \|x\| = 0 \Leftrightarrow x=0,\qquad \text{(N2)}\; \|\alpha x\| = |\alpha|\,\|x\|,\qquad \text{(N3)}\; \|x+y\| \le \|x\| + \|y\|.
    \]
    Das Paar \((V,\|\cdot\|)\) heißt normierter Raum.
\end{definition}

\begin{lemma}
    Sei \((V,\|\cdot\|)\) ein normierter Raum. Dann definiert
    \[
        d(x,y) := \|x-y\|
    \]
    eine Metrik auf \(V\).
\end{lemma}

\begin{lemma}
    Sei \((X,d)\) ein metrischer Raum und \(x,y,z \in X\). Dann gilt:
    \[
        |d(x,z) - d(y,z)| \le d(x,y).
    \]
\end{lemma}

\begin{definition}{Offene und abgeschlossene Kugeln}
    Sei \((X,d)\) ein metrischer Raum. Für \(x \in X\) und \(R>0\) definiert man:
    \[
        B_R(x) := \{y \in X : d(x,y) < R\},\qquad \overline{B}_R(x) := \{y \in X : d(x,y) \le R\}.
    \]
\end{definition}

\begin{definition}{Beschränkte Menge}
    Sei \((X,d)\) ein metrischer Raum. Eine Menge \(E \subset X\) heißt beschränkt, wenn es ein \(x \in X\) und \(R>0\) gibt mit
    \[
        E \subset B_R(x).
    \]
\end{definition}

\begin{definition}{Konvergenz von Folgen}
    Sei \((X,d)\) ein metrischer Raum und \((x_k)\) eine Folge in \(X\). Die Folge konvergiert gegen \(a \in X\), wenn gilt:
    \[
        \forall \varepsilon>0\;\exists k_0:\; d(x_k,a) < \varepsilon \text{ für alle } k \ge k_0.
    \]
    Man schreibt \(x_k \to a\).
\end{definition}

\begin{definition}{Häufungspunkt}
    Sei \((X,d)\) ein metrischer Raum und \((x_k)\) eine Folge in \(X\). Ein Punkt \(a \in X\) heißt Häufungspunkt der Folge, wenn eine Teilfolge \((x_{k_j})\) existiert mit
    \[
        x_{k_j} \to a.
    \]
\end{definition}

\begin{definition}{Konvergenz in \(\mathbb{R}^n\)}
    Sei \(x_k = (x_{k,1},\dots,x_{k,n}) \in \mathbb{R}^n\) und \(a = (a_1,\dots,a_n)\). Dann gilt:
    \[
        x_k \to a \Leftrightarrow \|x_k - a\|_\infty \to 0,\qquad \|x\|_\infty = \max_i |x_i|.
    \]

    Dies ist äquivalent zu:
    \[
        x_{k,i} \to a_i \text{ für alle } i.
    \]
\end{definition}

\begin{satz}{Bolzano-Weierstraß}
    Jede beschränkte Folge in \(\mathbb{R}^n\) besitzt eine konvergente Teilfolge.
\end{satz}

\begin{satz}{Äquivalenz der Normen}
    In \(\mathbb{R}^n\) sind alle Normen äquivalent. Zu jeder Norm \(\|\cdot\|\) existieren Konstanten \(m,M>0\) mit
    \[
        m\|x\|_\infty \le \|x\| \le M\|x\|_\infty \quad \text{für alle } x \in \mathbb{R}^n.
    \]
\end{satz}

\begin{definition}{Cauchy-Folge}
    Sei \((X,d)\) ein metrischer Raum. Eine Folge \((x_k)\) in \(X\) heißt Cauchy-Folge, falls für alle \(\varepsilon>0\) ein \(k^*\in\mathbb{N}\) existiert, sodass
    \[
        d(x_k,x_m) < \varepsilon \quad \text{für alle } k,m \ge k^*.
    \]
\end{definition}

\begin{satz}{Eigenschaften von Cauchy-Folgen}
    In einem metrischen Raum gilt:
    \begin{itemize}
        \item[(i)] Jede konvergente Folge ist eine Cauchy-Folge.
        \item[(ii)] Jede Cauchy-Folge ist beschränkt.
        \item[(iii)] Eine Cauchy-Folge hat höchstens einen Häufungspunkt. Sie konvergiert genau dann, wenn sie einen Häufungspunkt besitzt.
    \end{itemize}
\end{satz}

\begin{definition}{Vollständiger metrischer Raum}
    Ein metrischer Raum heißt vollständig, wenn jede Cauchy-Folge in ihm konvergiert.
\end{definition}

\begin{korollar}
    In einem vollständigen metrischen Raum ist eine Folge genau dann eine Cauchy-Folge, wenn sie konvergiert.
\end{korollar}

\begin{satz}{Vollständigkeit von \(\mathbb{R}^n\)}
    Der Raum \(\mathbb{R}^n\) ist vollständig. Als normierter Raum ist er daher ein Banachraum.
\end{satz}